% =============================================================================
% 020 / 068
% Status: raw
% =============================================================================
% Notes:
%
% Source question:
% 普通話能整理出文法嗎？普通話當然有文法呀！但是在這年頭，只能在正經的文章才能看見正常的文字了。網路上的用詞構句很散漫，新聞媒體則是精通語言，卻故意玩弄文字，而我竟習慣這些混亂的用法而不自知。怪我自己不愛讀書用腦哩⋯⋯不只你這樣批評過中文（普通話？）的現象，我聽過至少二十個立場相左、背景不同的人，都對此抱著相同的看法：現在的中文也好、普通話也罷，是很混亂的。
% =============================================================================

\chapter[普通話能整理出文法嗎？普通話當然有文法呀！但是在這年頭，只能在正經的文章才能看見正常的文字了。網路上的用詞構句很散漫，新聞媒體則是精通語言，卻故意玩弄文字，而我竟習慣這些混亂的用法而不自知。怪我自己不愛讀書用腦哩⋯⋯不只你這樣批評過中文（普通話？）的現象，我聽過至少二十個立場相左、背景不同的人，都對此抱著相同的看法：現在的中文也好、普通話也罷，是很混亂的。]{普通話能整理出文法嗎？普通話當然有文法呀！但是在這年頭，只能在正經的文章才能看見正常的文字了。網路上的用詞構句很散漫，新聞媒體則是精通語言，卻故意玩弄文字，而我竟習慣這些混亂的用法而不自知。怪我自己不愛讀書用腦哩⋯⋯不只你這樣批評過中文（普通話？）的現象，我聽過至少二十個立場相左、背景不同的人，都對此抱著相同的看法：現在的中文也好、普通話也罷，是很混亂的。}
\qaanswer{
普通話雖然語法混亂(這也能看出漢語使用者中其實是沒有什麼人文學菁英的，都近百年了，語法梳理還是如此混亂不堪，也沒有人在正經研究)，但因為使用人口的絕對巨大，很多使用方法甚至影響到了英語，大帝國的語言都有朝分析語發展的趨勢，所以她在全世界都是一門超強勢語言。但語言的強勢和語言的美(既精確又經濟)不能劃等號。
}

